% Section 2, part 7 of 8: the model cohomological field theory; its homogeneity; the dilaton lemma.

\begin{defn}[The model cohomological field theory]
\label{def:set-modcohft}
Let $\varsigma$ be a formal fifth root of $v$, as in \cref{prop:set-modrows}. The \emph{model TFT} $\TFT^{\mathrm m}_v$ is the TFT on the colour space, with the orthonormal basis $\tilde e_i$ ($i\in\bbZ/5$), \emph{defined} by $\hat\Delta_i^{1/2}\coloneqq\zeta^{4i}\varsigma$:
\[
\TFT^{\mathrm m}_{v;h,n}(\tilde e_i,\dots,\tilde e_i)\coloneqq(\zeta^{4i}\varsigma)^{2h-2+n},\qquad\TFT^{\mathrm m}_{v;h,n}=0\text{ on mixed colours},
\]
with the unit $\sum_i\hat\Delta_i^{-1/2}\tilde e_i=\varsigma^{-1}\sum_i\zeta^i\tilde e_i$. Let $\Psi^{\mathrm m}_v$ be the map from the colour space to $\bbH$ with the lowered entries $\eta(H^b,\Psi^{\mathrm m}_v\tilde e_i)\coloneqq\varsigma^{\nu_b-1}\varepsilon_b\zeta^{i(b+1)}=\varsigma^{\nu_b-1}(\Psi_0)_{bi}$: the value at $z=0$ of the colour columns of \cref{def:set-M0} with the exponentials stripped (the model columns $\Qr^{\mathrm m}$ of \cref{thm:gau-R0} at $z=0$). It is an isometry: on the pairs $a+b\equiv3$ the powers of $\varsigma$ cancel because $\nu_a+\nu_b=2$, so $(\Psi^{\mathrm m}_v)^T\eta^{-1}\Psi^{\mathrm m}_v=\Psi_0^T\eta^{-1}\Psi_0=1$ (\cref{def:set-M0}). The \emph{model cohomological field theory} on $(\bbH,\eta)$ is
\[
\Omega^{\mathrm{mod}}_v\coloneqq(\Psi^{\mathrm m}_v)_*\bigl(\Mo(z/\varsigma^4).\TFT^{\mathrm m}_v\bigr),\qquad(A_*\Omega)_{h,n}\coloneqq\Omega_{h,n}\circ(A^{-1})^{\otimes n},
\]
where the action is the unit-preserving action of \cite[Def.~2.13]{PPZ15} (with \cite[Def.~2.5]{PPZ15}), and $\Mo(z/\varsigma^4)=1+O(z)$ is symplectic by \cref{lem:set-M0sympl}. Its ancestor correlators are
\[
\langle\psi^{k_1}H^{a_1}\cdots\psi^{k_n}H^{a_n}\rangle_h(v)\coloneqq\int_{\overline{\mathcal M}_{h,n}}\Omega^{\mathrm{mod}}_{v;h,n}(H^{a_1},\dots,H^{a_n})\prod_j\psi_j^{k_j}\qquad(2h-2+n>0),
\]
and its dilaton-shifted potentials $\overline F_h(q)$, with the base point $q=-z1$, are those of \cref{def:set-graphsum}. The \emph{free energies} are $\Fm_h\coloneqq\langle\ \rangle_h(v)$ ($h\geq2$); that $\Fm_h=c_hv^{2-2h}$ is \cref{thm:mod-gap}. The identification of $\Omega^{\mathrm{mod}}_v$ with the theory of a Hurwitz Frobenius manifold is not used (\cref{sec:sgn}).
\end{defn}

\begin{lem}[Homogeneity of the model theory]
\label{lem:set-modhom}
\begin{itemize}
\item[(a)] For every cohomological degree $D$,
\[
\Omega^{\mathrm{mod}}_{v;h,n}(H^{b_1},\dots,H^{b_n})_{[D]}=\varsigma^{\,2h-2+\sum_j\nu_{b_j}-4D}\,\Omega^{\mathrm{mod}}_{1;h,n}(H^{b_1},\dots,H^{b_n})_{[D]},
\]
where ${}_{[D]}$ is the part of degree $D$ and $\Omega^{\mathrm{mod}}_1$ is obtained by putting $\varsigma=1$.
\item[(b)] Every ancestor correlator is $\langle\prod_j\psi^{k_j}H^{b_j}\rangle_h(v)=C\,v^e$ with $C\coloneqq\langle\prod_j\psi^{k_j}H^{b_j}\rangle_h(1)$ and
\[
e\coloneqq2-2h+\sum_j\wt(k_j,b_j),\qquad\wt(k,b)\coloneqq\frac{4k+\nu_b-4}5,\qquad v^e\coloneqq\varsigma^{5e}.
\]
\item[(c)] $C=0$ unless $e\in\bbZ$, and $C\in\bbQ$. A correlator is \emph{regular} if $e\geq0$; then it lies in $\bbQ[v]$.
\end{itemize}
\end{lem}

\begin{proof}
(a) First for $\Omega'\coloneqq\Mo(z/\varsigma^4).\TFT^{\mathrm m}_v$ in the basis $\tilde e_i$, graph by graph in \cite[Def.~2.2]{PPZ15} with the translation of \cite[Def.~2.13]{PPZ15}, $\ttr(z)=z1'-z\Mo(z/\varsigma^4)^{-1}1'$, where $1'=\varsigma^{-1}\sum_i\zeta^i\tilde e_i$ is the unit of $\TFT^{\mathrm m}_v$. Every ingredient is free of $\varsigma$ except the following. A leg monomial $\psi^k$ of $\Mo^{-1}(\psi/\varsigma^4)$ carries $\varsigma^{-4k}$. An edge $\bigl(1-\Mo^{-1}(\psi'/\varsigma^4)\Mo^{-1}(\psi''/\varsigma^4)^T\bigr)/(\psi'+\psi'')$ is $\varsigma^{-4}$ times a series in $\psi'/\varsigma^4$, $\psi''/\varsigma^4$, so its monomial $\psi'^a\psi''^b$, of degree $a+b+1$ after the push-forward along the gluing map, carries $\varsigma^{-4(a+b+1)}$. The coefficient of $z^m$ in $\ttr(z)$ is proportional to $\varsigma^{-1}\varsigma^{-4(m-1)}$; a translation point with $\psi^m$ contributes the degree $m-1$ after the push-forward along the forgetful map, and it raises the valence of its vertex by one. A vertex of genus $g_v$ with $n_v$ half-edges and legs and $m_v$ translation points carries $\varsigma^{2g_v-2+n_v+m_v}$, from $\hat\Delta_i^{1/2}=\zeta^{4i}\varsigma$. Since $\sum_v(2g_v-2+n_v)=2h-2+n$ for a stable graph of genus $h$ with $n$ legs, and each translation point contributes $1+(-1-4(m-1))=-4(m-1)$ to the exponent of $\varsigma$ and $m-1$ to the degree, every term of degree $D$ carries $\varsigma^{2h-2+n-4D}$. Finally $(\Psi^{\mathrm m}_v)^{-1}H^b=\sum_i\eta(H^b,\Psi^{\mathrm m}_v\tilde e_i)\tilde e_i=\varsigma^{\nu_b-1}\sum_i\varepsilon_b\zeta^{i(b+1)}\tilde e_i$, as $\Psi^{\mathrm m}_v$ is an isometry, and the $n$ insertions contribute $\varsigma^{\sum_j(\nu_{b_j}-1)}$. This gives (a).

(b) Integrating against $\prod_j\psi_j^{k_j}$ selects the degree $D=3h-3+n-\sum_jk_j$, and the exponent of $\varsigma$ becomes $2h-2+\sum_j\nu_{b_j}-4(3h-3+n-\sum_jk_j)=5\bigl(2-2h+\sum_j(4k_j+\nu_{b_j}-4)/5\bigr)=5e$.

(c) Replace $\varsigma$ by $\zeta^{-1}\varsigma$. By step~(1) of the proof of \cref{lem:set-M0sympl}, this replaces, for every $i$, the colour-$i$ column of \cref{def:set-M0} by the colour-$(i+1)$ column; so it replaces $\Psi^{\mathrm m}_v$ by $\Psi^{\mathrm m}_v$ composed with the cyclic permutation $\tilde e_i\mapsto\tilde e_{i+1}$ of the colour space and $\Mo(t)$ by its conjugate by that permutation, and it replaces $\hat\Delta_i^{1/2}=\zeta^{4i}\varsigma$ by $\zeta^{4i-1}\varsigma=\hat\Delta_{i+1}^{1/2}$. By the naturality $A_*(R.\Omega)=(ARA^{-1}).(A_*\Omega)$ of \cref{def:set-graphsum}, applied to the permutation, an isometry of the colour space, $\Omega^{\mathrm{mod}}_v$ in the flat basis does not depend on the choice of the fifth root. By (b), $(\zeta^{-5e}-1)C=0$, so $C=0$ unless $e\in\bbZ$. In the same way an automorphism $\zeta\mapsto\zeta^k$ of $\bbQ(\zeta)$ ($k\in(\bbZ/5)^\times$) fixes $\varsigma$ and the series $\Fs_b$, which lie in $\bbQ\fps t$ because the colour-$0$ T-system has rational coefficients; it maps the colour-$i$ column to the colour-$ki$ column and $\hat\Delta_i^{1/2}$ to $\hat\Delta_{ki}^{1/2}$. By the same naturality, applied to the permutation $\tilde e_i\mapsto\tilde e_{ki}$, $\Omega^{\mathrm{mod}}_v$ in the flat basis is invariant under the Galois group of $\bbQ(\zeta)$; hence its correlators have rational coefficients, and $C\in\bbQ$. If $e\geq0$, $Cv^e\in\bbQ[v]$.
\end{proof}

\begin{lem}[The dilaton equation and the homogeneity of the potentials]
\label{lem:set-dilaton}
\begin{enumerate}
\item Let $\Omega$ be a cohomological field theory with flat unit $1$ \cite[Def.~0.3~(i)--(iii)]{PPZ15}, with correlators as in \cref{def:set-graphsum}. For $2h-2+n>0$,
\[
\langle\psi1\cdot\psi^{k_1}a_1\cdots\psi^{k_n}a_n\rangle_h=(2h-2+n)\,\langle\psi^{k_1}a_1\cdots\psi^{k_n}a_n\rangle_h .
\]
\item $\Omega^{\mathrm{mod}}_v$ has the flat unit $H^0=1$.
\item Let $\mathsf R$ be a complete local $\bbQ$-algebra with maximal ideal $\mathfrak m$ (for instance $\bbQ\fps u$) and $\varkappa\in1+\mathfrak m$, or let $\mathsf R=\bbQ(\zeta)\lser s$ with the $s$-adic topology and $\varkappa\in1+u\bbQ\fps u$ ($u=s^5$); let $\Omega$ be a cohomological field theory with flat unit $1$ on a free $\mathsf R$-module $V$ of finite rank. Regard $\overline F_h(q)=\sum_n\frac1{n!}\langle(q+z1)^n\rangle_h$ as a formal function of $y\coloneqq q+z1\in V\fps z$, $y$ being a formal variable. Then, for $2h-2+n>0$,
\[
\dd^n\overline F_h(q/\varkappa)=\varkappa^{2h-2+n}\,\dd^n\overline F_h(q),
\]
and, for $n=0$ and $h\neq1$, $\overline F_h(\varkappa q)=\varkappa^{2-2h}\,\overline F_h(q)$.
\end{enumerate}
\end{lem}

\begin{proof}
(i) Let $\pi\colon\overline{\mathcal M}_{h,n+1}\to\overline{\mathcal M}_{h,n}$ forget the last point; by the symmetry of the classes we may put the insertion $\psi1$ there. By \cite[Def.~0.3~(iii)]{PPZ15}, $\Omega_{h,n+1}(a_1,\dots,a_n,1)=\pi^*\Omega_{h,n}(a_1,\dots,a_n)$. By \cite[(1.7)]{AC96}, $\pi_*\bigl(\psi_1^{k_1}\cdots\psi_n^{k_n}\psi_{n+1}^{a+1}\bigr)=\psi_1^{k_1}\cdots\psi_n^{k_n}\kappa_a$, the classes $\psi_j$ on the left being those of $\overline{\mathcal M}_{h,n+1}$, and $\kappa_0=2h-2+n$ \cite[p.~5]{AC96}. With $a=0$ and the projection formula, $\int\pi^*\Omega\cdot\psi_{n+1}\prod_j\psi_j^{k_j}=\int\Omega\cdot\pi_*\bigl(\psi_{n+1}\prod_j\psi_j^{k_j}\bigr)=(2h-2+n)\int\Omega\prod_j\psi_j^{k_j}$.

(ii) By \cite[Prop.~2.12]{PPZ15} with \cite[Def.~2.13]{PPZ15}, $\Mo(z/\varsigma^4).\TFT^{\mathrm m}_v$ is a cohomological field theory with the unit $\varsigma^{-1}\sum_i\zeta^i\tilde e_i$ of $\TFT^{\mathrm m}_v$, and the push-forward by an isometry preserves the axioms. The lowered components of the image of this unit under $\Psi^{\mathrm m}_v$ are $\sum_i\zeta^i\varsigma^{-1}\varsigma^{\nu_b-1}\varepsilon_b\zeta^{i(b+1)}=\varepsilon_b\varsigma^{\nu_b-2}\sum_i\zeta^{i(b+2)}=5\,\delta_{b3}$, since $\varepsilon_3=1$ and $\nu_3=2$; these are the lowered components $\eta(H^b,1)$ of $H^0=1$.

(iii) Write $q/\varkappa+z1=y/\varkappa+\epsilon z1$ with $\epsilon\coloneqq1-\varkappa^{-1}$, which lies in $\mathfrak m$, respectively in $u\bbQ\fps u$, so that $\epsilon^b\to0$ in $\mathsf R$; and $x=(x_1,\dots,x_n)$. By multilinearity, $\frac1{m!}\langle x\,(y/\varkappa+\epsilon z1)^m\rangle_h=\sum_{a+b=m}\frac{\varkappa^{-a}\epsilon^b}{a!\,b!}\langle x\,y^a(z1)^b\rangle_h$. Applying (i) $b$ times, which is allowed because $2h-2+n+a>0$ at every step, $\langle x\,y^a(z1)^b\rangle_h=(N)_b\langle x\,y^a\rangle_h$ with $N\coloneqq2h-2+n+a$ and $(N)_b\coloneqq N(N+1)\cdots(N+b-1)$. Hence
\[
\dd^n\overline F_h(q/\varkappa)[x]=\sum_a\frac{\varkappa^{-a}}{a!}\sum_b\frac{\epsilon^b}{b!}(N)_b\,\langle x\,y^a\rangle_h=\sum_a\frac{\varkappa^{-a}}{a!}(1-\epsilon)^{-N}\langle x\,y^a\rangle_h=\varkappa^{2h-2+n}\,\dd^n\overline F_h(q)[x],
\]
by the binomial series, which converges in $\mathsf R$ because $\epsilon^b\to0$, and $(1-\epsilon)^{-1}=\varkappa$. The identity holds coefficientwise in $y$, as each $\langle x\,y^a(z1)^b\rangle_h$ is a homogeneous polynomial of degree $a$ in $y$. The statement for $n=0$ follows in the same way, with $\varkappa^{-1}$, which satisfies the same hypothesis, in place of $\varkappa$ and $N=2h-2+a$; for $h=0$ the terms with fewer than three insertions $y$ vanish on both sides, because $\psi=0$ on $\overline{\mathcal M}_{0,3}$.
\end{proof}

The hypothesis on $\varkappa$ is used: for $\mathsf R=\bbQ\fps u$, and for $\mathsf R=\bbQ(\zeta)\lser s$, it says $\varkappa(0)=1$, and only then are the identities of (iii) identities of Taylor expansions at the base point $-z1$; for a constant $\varkappa\neq1$ the series in $\epsilon$ does not converge. In the applications (\cref{lem:tra-F}, \cref{lem:tra-conv}~(f), \cref{lem:flh-dil}, the last over both rings) the scalar $\varkappa$ is the dilaton factor $d\in\bbQ\fps u$ of the canonical gauge, and $d(0)=1$ by \cref{lem:gau-diag}.

% Section 2, part 8 of 8: the hatted theory, the outer graph sum, polar parts by reversion; Property N(g), the constants c_g.

%-------------------------------------------------------------------------------
\subsection{The hatted theory and the outer graph sum}
\label{sec:set-hat}

In this subsection a conifold frame is fixed, and the rows are expanded in it (\cref{def:set-confexp}; they lie in $\Nring\otimes\bbQ(\zeta)$ by \cref{cor:gau-rowsNN}): $\Qr^{(i)}_a(z)$ is a series in $z$ whose coefficients are Laurent series in $s$, and so is every series below.

\begin{defn}[The hatted theory]
\label{def:set-hat}
The \emph{hatted classes} are $\hat H^a\coloneqq s^{1-\nu_a}H^a$, the hatted variable is $\zh=z/s^4$, and the \emph{hatted rows} are the lowered components with respect to the $\hat H^a$, as series in $\zh$:
\[
\hat\Qr^{(i)}_a(\zh)\coloneqq s^{1-\nu_a}\,\Qr^{(i)}_a(s^4\zh).
\]
The pairing is preserved: $\eta(\hat H^a,\hat H^b)=s^{2-\nu_a-\nu_b}\eta(H^a,H^b)=\eta(H^a,H^b)$, since $\nu_a+\nu_b=2$ on the pairs $a+b\equiv3$. The hatted unit of the state space is $\hat1\coloneqq\hat H^0=s1$. The element $\zeta^{-4i}\Cexp_{\alpha,\beta}(\Delta_i^{1/2})=s\varpi_0$ does not depend on $i$; put $\hat x\coloneqq\varpi_0=1+\alpha u+O(u^2)$. The \emph{hatted TFT} is $\hat\TFT\coloneqq s^{-1}.\TFT$, i.e.\ $\hat\TFT_{g,n}\coloneqq s^{-(2g-2+n)}\TFT_{g,n}$ (the rescaling of \cref{def:flh-outer}~(iv)): on the colour space $\hat\TFT_{g,n}(\tilde e_i,\dots,\tilde e_i)=(\zeta^{4i}\hat x)^{2g-2+n}$, and its unit is $s1_{\TFT}$. That the hatted rows are power series in $u$ is not assumed here; it is proved in \cref{cor:mai-natfact}.
\end{defn}

\begin{defn}[The outer graph sum]
\label{def:set-outer}
Let $M\coloneqq\Mo(\tau(u)\zh)$, taken exactly in $u$: with $\zh=z/s^4$ it is a power series in $z$ with coefficients in $\bbQ(\zeta)\lser s$, with $M(0)=1$ and $M(-z)^TM(z)=1$ (\cref{def:set-M0}, \cref{lem:set-M0sympl}). Let $\Qr$ be the matrix of the lowered rows of \cref{def:set-rows} expanded in the frame, and
\[
R'\coloneqq\Qr\,M^{-1},
\]
a series of maps from the colour space (orthonormal basis $\tilde e_i$) to $\bbH$. Its constant term is $R'(0)=\Qr(0)$, the expansion of the matrix $\Psi_{\mathrm Q}$ with the entries $\Qr^{(i)}_{a,0}$, the representatives in $\Nring\otimes\bbQ(\zeta)$ of $\eta(H^a,\tilde e_i)$. By the explicit form of $\Psi_{\mathrm Q}$ in \cref{lem:mai-Psi}~(i), $(\Psi_{\mathrm Q})_{ai}=\zeta^{i(a+1)}(\Psi_{\mathrm Q})_{a0}$ with $(\Psi_{\mathrm Q})_{30}=C_0/L$, and $\Psi_{\mathrm Q}^T\eta^{-1}\Psi_{\mathrm Q}=1$; with $\Delta_i^{-1/2}=\zeta^iL/C_0$ this gives $\sum_i\Delta_i^{-1/2}(\Psi_{\mathrm Q})_{ai}=(L/C_0)(\Psi_{\mathrm Q})_{a0}\sum_i\zeta^{i(a+2)}=5\,\delta_{a3}$. These are identities in $\Nring\otimes\bbQ(\zeta)$, so $R'(0)$ is an isometry which maps $1_{\TFT}$ to $1$. Its matrix is not the identity, and $R'$ satisfies the condition (unit) of \cref{def:flh-outer} with $\varkappa=1$, a condition on $R'(0)$ only. The \emph{outer graph sum} has the vertices $(M.\TFT)_{g,n}$, the classes of the cohomological field theory $M.\TFT$ on the colour space \cite[Def.~2.13]{PPZ15}, and the edges and leaves of $R'$: with $N(\psi,\psi')\coloneqq1-\bigl(R'^{-1}(\psi)\otimes R'^{-1}(\psi')\bigr)\eta^{-1}$, where $1=\sum_i\tilde e_i\otimes\tilde e_i$ is the inverse bivector of the colour space,
\[
E_{R'}(\psi,\psi')\coloneqq\frac{N(\psi,\psi')-\frac12N(-\psi',\psi')-\frac12N(\psi,-\psi)}{\psi+\psi'},\qquad q_{R'}(z)\coloneqq R'^{-1}(z)(-z1),
\]
with the leaf insertion $q_{R'}+z1_{\TFT}\in z^2(\text{colour space})\fps z$. The numerator vanishes at $\psi'=-\psi$, so $E_{R'}$ is a power series in $\psi,\psi'$; it is symmetric under the exchange of its two slots together with $\psi$ and $\psi'$, and its coefficient of $\psi^a\psi'^b$ involves $R'^{-1}$ only through $z^{a+b+1}$. No property of $R'$ other than the invertibility of $R'(0)$ is used. If $R'$ is symplectic, $R'(-z)^T\eta^{-1}R'(z)=1$, then $N(-\psi',\psi')=N(\psi,-\psi)=0$, $E_{R'}$ is the edge of \cite[Def.~2.2]{PPZ15} \cite[Rem.~2.1, (6)]{PPZ15}, and the outer graph sum is the leg-free part of $R'.(M.\TFT)$ in the sense of \cref{def:flh-outer}. For a leg-free stable graph $\Gamma$ of genus $g\geq2$ the \emph{outer term} is
\[
T_\Gamma\coloneqq\mathrm{Cont}_\Gamma\bigl[E_{R'},q_{R'};M.\TFT\bigr]\in\bbQ(\zeta)\lser s
\]
(\cref{def:set-graphsum}); the vertex tensors are the full tensors of $M.\TFT$, so $T_\Gamma$ contains the sum over all colourings of $\Gamma$. The \emph{vertex graph} is the stable graph with one vertex, of genus $g$, and no edge; its outer term contains all the leaf insertions of $R'$. The \emph{hatted term} $G_\Gamma$ is the same contraction formed with the hatted data of \cref{def:set-hat}: with $\hat M(\zh)\coloneqq\Mo(\tau\zh)$ and $\hat R'(\zh)\coloneqq\hat\Qr(\zh)\hat M(\zh)^{-1}$,
\[
G_\Gamma\coloneqq\mathrm{Cont}_\Gamma\bigl[E_{\hat R'},q_{\hat R'};\hat M.\hat\TFT\bigr],
\]
the edge and the leaf point being formed in the variable $\zh$ and with the unit $\hat1$. By \cref{lem:flh-hat}, $G_\Gamma=u^{2g-2}T_\Gamma$. The outer term $T_\Gamma$ is \emph{regular at the conifold point} if it lies in $\bbQ\fps u$.

\emph{The symplecticity of $R'$.} As series in $q$ the rows are symplectic (\cref{def:set-rows}~(a)): $\Qr(-z)^T\eta^{-1}\Qr(z)=1$ over $\bbQ(\zeta)\fps q$. This identity is not transferred to the conifold expansion: that would require it in $\Nring\otimes\bbQ(\zeta)$, where it is not proved, and the injectivity of the evaluation $\ev$ of \cref{def:set-confexp} is not claimed. Instead, by \cref{cor:gau-colours} and \cref{lem:flh-inputs}~(i), the expanded rows are $\Qr=\Pc\Psi^{\mathrm m}_{v(u)}M\Dn^{-1}$, with $\Dn=\diag(\Dn_i)$ the normalisation series of \cref{thm:gau-R0} and \cref{cor:gau-colours} and $\Psi^{\mathrm m}_{v(u)}$ the isometry of \cref{def:set-modcohft}. With $\Pc(-z)^T\eta^{-1}\Pc(z)=\eta^{-1}$ (\cref{prop:gau-symp}; $\eta^{-1}=\eta/25$), $(\Psi^{\mathrm m}_{v(u)})^T\eta^{-1}\Psi^{\mathrm m}_{v(u)}=1$ and $M(-z)^TM(z)=M(z)M(-z)^T=1$ (\cref{lem:set-M0sympl}) this gives
\[
R'(-z)^T\eta^{-1}R'(z)=M(z)\,\diag\bigl((\Dn_i(z)\Dn_i(-z))^{-1}\bigr)_i\,M(-z)^T .
\]
So $R'$ is symplectic modulo $z^{3g-2}$ when $\Fg{g}$ holds (\cref{def:gau-Fg}), and it is symplectic because $\Dn=1$ (\cref{thm:gam-D}). The proof of \cref{thm:gam-D} uses neither $R'$ nor $T_\Gamma$, so there is no circularity, and no statement of this section uses the symplecticity of $R'$.
\end{defn}

\begin{rmk}[Polar parts by reversion]
\label{rmk:set-reversion}
Let $g\geq2$, $G\in\bbQ\fps u$ and $T\coloneqq u^{2-2g}G$; for instance $T=T_\Gamma$ and $G=G_\Gamma$, once the hatted data are known to be power series in $u$ (\cref{cor:mai-natfact}). In a conifold frame $\tc=u(1+O(u))\in\bbQ\fps u$, so $\bbQ\fps u=\bbQ\fps{\tc}$ and $u=\tc(1+O(\tc))$. For $m\in\bbZ$, $u^m=\tc^m(u/\tc)^m$ with $(u/\tc)^m\in1+\tc\bbQ\fps{\tc}$; hence
\[
T=\sum_{k=2-2g}^{-1}b_k\,\tc^k+(\text{an element of }\bbQ\fps u),
\]
where $(b_{2-2g},\dots,b_{-1})$ is the image of $(a_0,\dots,a_{2g-3})$, $a_k\coloneqq\coef{u^k}G$, under a lower triangular matrix with unit diagonal that depends only on $\tc$. So $T$ is regular if and only if $a_0=\dots=a_{2g-3}=0$, and $T\in c\,\tc^{2-2g}+\bbQ\fps u$ if and only if $a_0=c$ and $2g-3$ further affine conditions on $a_1,\dots,a_{2g-3}$ hold. For $g=2$, with $\gamma_t=\coef{u^2}\tc$ (\cref{def:set-frames}), $T=a_0\tc^{-2}+(a_1+2\gamma_ta_0)\tc^{-1}+\bbQ\fps u$.
\end{rmk}

%-------------------------------------------------------------------------------
\subsection{Property N(\texorpdfstring{$g$}{g})}
\label{sec:set-N}

\begin{defn}[Property N($g$)]
\label{def:set-N}
Let $g\geq2$ and
\begin{equation}
\label{eq:set-cg}
c_g\coloneqq\frac{(-1)^{g-1}B_{2g}\,5^{1-g}}{2g(2g-2)}=\frac{|B_{2g}|}{2g(2g-2)\,5^{g-1}} .
\end{equation}
\emph{Property N($g$)} holds if, for every conifold frame of \cref{def:set-frames}, the outer terms of \cref{def:set-outer}, with $M=\Mo(\tau(u)\zh)$ exactly in $u$, satisfy:
\begin{enumerate}
\item $T_\Gamma\in\bbQ\fps u$ for every leg-free stable graph $\Gamma$ of genus $g$ with at least one edge;
\item $T_\Gamma\in c_g\,\tc^{2-2g}+\bbQ\fps u$ for the vertex graph.
\end{enumerate}
Property N($g$) is a statement about every graph, not only about the sum over the graphs. It is the statement of \cref{thm:mai-main} for one $g$.
\end{defn}

\begin{rmk}[The constants $c_g$]
\label{rmk:set-cg}
As $(-1)^{g-1}B_{2g}=|B_{2g}|$, the two forms in \eqref{eq:set-cg} agree and $c_g>0$. With $B_4=-\frac1{30}$, $B_6=\frac1{42}$, $B_8=-\frac1{30}$ and $B_{10}=\frac5{66}$,
\[
c_2=\frac1{1200},\qquad c_3=\frac1{25200},\qquad c_4=\frac1{180000},\qquad c_5=\frac1{660000}.
\]
In the model, $c_h$ is the constant of the free energies $\Fm_h=c_hv^{2-2h}$ (\cref{thm:mod-gap}), and
\[
c_h=(-5)^{1-h}\,\frac{B_{2h}}{2h(2h-2)}
\]
is $(-5)^{1-h}$ times the coefficient of the free energy $B_{2h}v^{2-2h}/(2h(2h-2))$ of topological recursion on the model curve (\cref{cor:mod-trF}). The constant refers to the normalisation $\varpi_1=u+O(u^2)$ of the vanishing period: for the vanishing period $k\varpi_1$, $k\in\bbQ^\times$, the flat coordinate is $k\tc$, and $c_g\tc^{2-2g}=(k^{2g-2}c_g)(k\tc)^{2-2g}$.
\end{rmk}

